\documentclass[14pt,a4paper]{extarticle}
\usepackage[margin=2cm]{geometry}
\usepackage{amsmath,amssymb,mathtools,bm}
\usepackage{graphicx}
\usepackage{float}
\usepackage{booktabs}
\usepackage[T1]{fontenc}
\usepackage{lmodern}
\usepackage{xcolor}
\usepackage[colorlinks=true,linkcolor=blue!55!black]{hyperref}
\numberwithin{equation}{section}
\setlength{\parindent}{0pt}
\setlength{\parskip}{0.55em}
\linespread{1.2}
\newcommand{\M}{\mathcal M}
\newcommand{\ii}{\mathrm i}
\newcommand{\slsh}[1]{\not\!#1}
\title{SIDIS}
\author{Caiyi, Jiang}
\date{}
\begin{document}
\maketitle
\setcounter{tocdepth}{2}
\tableofcontents
\clearpage

\section{Factorization notes}
\label{sec:factorization-notes}
Based on Sec.~2.4 of Zhite Yu's thesis,
\href{https://arxiv.org/abs/2308.13080}{arXiv:2308.13080v1}
\subsection{Basic concepts}
\label{sec:factorization-foundations}
\subsubsection{Sudakov form factor}
The form factor describes QCD corrections to the quark electromagnetic vertex:
\begin{equation}
 \langle q(k')|j_{\rm em}^\mu|q(k)\rangle
 =e_a\bar u(k')\gamma^\mu u(k)F(Q^2),\qquad Q^2=-(k'-k)^2.
\end{equation}
Here $e_a$ is the fractional charge and $F=1$ at tree level.

\subsubsection{Pinch surfaces}
A pinch occurs when poles trap a loop integration contour. For example,
\begin{equation}
 I(a,b)=\int_{-\infty}^{\infty}
 \frac{dk^0}{(k^0-a+\ii0)(k^0-b-\ii0)}
\end{equation}
has poles below and above the real axis. They pinch the contour as $a\to b$.
In several dimensions, the trapped momentum configurations form a pinch
surface. \textbf{Nearby regions can give long-distance contributions.}
\subsubsection{Collinear and hard subgraphs}
A subgraph is a connected set of lines and vertices. In this example,
\begin{equation}
 \underbrace{C_q}_{+z\text{ collinear}}\longleftrightarrow
 \underbrace{H}_{\text{hard}}\longleftrightarrow
 \underbrace{C_{\bar q}}_{-z\text{ collinear}}.
\end{equation}

In $C_q$ and $C_{\bar q}$, momenta point nearly in the same direction
and have small virtualities. In $H$, internal virtualities are of order
$Q^2$. 
\textbf{Each $C$ joins $H$ through one fermion and possibly many gluons.}

All collinear propagators, including those next to $H$, belong to $C$.
For example,
\begin{equation}
 C_{q,\alpha}(k)=B_\beta(k)
 \left[\frac{\ii(\slsh k+m)}{k^2-m^2+\ii0}\right]_{\beta\alpha},
 \qquad
 \mathcal A=\int\frac{d^4k}{(2\pi)^4}C_{q,\alpha}(k)H_\alpha(k).
\end{equation}
$B_\beta(k)$ is the rest of the collinear quark subgraph before attaching the propagator that connects it to the hard subgraph.
Thus $H$ has no extra connecting quark propagator. 
For gluons, the scalar propagator factor belongs to $C$.

\subsubsection{Light-cone components and virtuality}
Use $g=\operatorname{diag}(1,-1,-1,-1)$ and $c=1$. Then
\begin{align}
 k^\mu&=(k^0,k^1,k^2,k^3)=(E,k_x,k_y,k_z),\\
 k^\pm&=\frac{k^0\pm k^3}{\sqrt2},\qquad
 k^0=\frac{k^++k^-}{\sqrt2},\quad k^3=\frac{k^+-k^-}{\sqrt2},\\
 k^2&=2k^+k^--\bm k_T^2,\qquad
 k\cdot l=k^+l^-+k^-l^+-\bm k_T\cdot\bm l_T.
\end{align}
An on-shell particle has $k^2=m^2$; its off-shellness is $k^2-m^2$.
For a massless line, the virtuality is $k^2$; its size is $|k^2|$.
The SIDIS photon has $q^2=-Q^2$.

\textbf{1. Large component: motion near $+z$.}
\begin{equation}
 k^0\simeq k^3>0
 \quad\Rightarrow\quad
 k^+=\frac{k^0+k^3}{\sqrt2}\sim Q,
 \qquad k^-=\frac{k^0-k^3}{\sqrt2}\ll Q.
\end{equation}

\textbf{2. Small angle: small transverse momentum.}
\begin{equation}
 |\bm k_T|\sim\lambda\ll Q,
 \qquad \theta\simeq\frac{|\bm k_T|}{|k^3|}
 \sim\frac{\lambda}{Q}\ll1.
\end{equation}

\textbf{3. Small virtuality: solve for the minus component.}
\begin{align}
 k^2&=2k^+k^--\bm k_T^2=O(\lambda^2),\\
 2k^+k^-&=\underbrace{k^2}_{O(\lambda^2)}
          +\underbrace{\bm k_T^2}_{O(\lambda^2)},\\
 k^-&=\frac{k^2+\bm k_T^2}{2k^+}
 \sim\frac{\lambda^2}{Q}.
\end{align}
This gives $k^+\gg|\bm k_T|\gg k^-$ for typical collinear momenta.

\textbf{4. Motion near $-z$: exchange plus and minus.}
\begin{align}
 l^0\simeq-l^3>0
 &\quad\Rightarrow\quad l^-=\frac{l^0-l^3}{\sqrt2}\sim Q,\\
 |\bm l_T|\sim\lambda,\quad l^2=O(\lambda^2)
 &\quad\Rightarrow\quad
 l^+=\frac{l^2+\bm l_T^2}{2l^-}\sim\frac{\lambda^2}{Q}.
\end{align}
Thus, in the order $(+,-,T)$, source Eq.~(2.56) gives
\begin{equation}
 \boxed{k_i\sim(Q,\lambda^2/Q,\lambda),\qquad
 l_j\sim(\lambda^2/Q,Q,\lambda).}
 \label{eq:notes-collinear-scaling}
\end{equation}
The order is $(+,-,T)$; $i=q,1,\ldots,n$ and $j=\bar q,1,\ldots,m$.
$\lambda$ has units of momentum, and $\lambda/Q\ll1$.


\subsubsection{Why extra gluons matter}
More gluons mean higher orders in the coupling, but not always extra
powers of $\lambda/Q$. \textbf{Longitudinal gluon contractions can remain leading.}
These are internal gauge-field components, not extra physical gluon helicities.


\subsubsection{Ward identity}
\label{sec:notes-ward-identity}
A Ward identity relates a gauge-boson momentum contraction to differences
of inverse propagators. Start with the tree-level vector vertex.

\textbf{1. Contract the vertex with the momentum.}
For a fermion that absorbs momentum $k$, $p'=p+k$. Then
\begin{equation}
 k_\mu\gamma^\mu=\slsh k
 =(\slsh p'-m)-(\slsh p-m).
\end{equation}
For on-shell external spinors, the Dirac equations give
\begin{align}
 \bar u(p')(\slsh p'-m)&=0,\qquad
 (\slsh p-m)u(p)=0,\\
 k_\mu J^\mu
 &=\bar u(p')\bigl[(\slsh p'-m)-(\slsh p-m)\bigr]u(p)
 =0,
 \qquad J^\mu=\bar u(p')\gamma^\mu u(p).
\end{align}
Thus this current is unchanged in the amplitude when the photon
polarization is shifted by a multiple of its momentum.

\textbf{2. Insert the contracted vertex on an internal line.}
To show the algebra, omit factors of $\ii$ and pole prescriptions and set
\begin{equation}
 S_0(p)=(\slsh p-m)^{-1},\qquad
 \slsh k=S_0^{-1}(p+k)-S_0^{-1}(p).
\end{equation}
Multiplication on both sides gives
\begin{align}
 S_0(p+k)\slsh k S_0(p)
 &=S_0(p+k)
 \bigl[S_0^{-1}(p+k)-S_0^{-1}(p)\bigr]S_0(p)\\
 &=\boxed{S_0(p)-S_0(p+k)}.
\end{align}
The momentum insertion removes one adjacent propagator in each term.
When all allowed attachments are summed, such differences cancel between
neighbouring terms and leave endpoint terms.

\textbf{3. Use it in factorization.}
The projections below turn gluon connections into
\begin{equation}
 \hat k^\nu H_\nu
 \qquad\text{or}\qquad
 \hat k_s^\mu C_\mu.
\end{equation}
These are momentum contractions of the same type. In QCD the full
identity also involves color and gluon self-interactions, so the required
diagram sums must be included. The remaining attachments form Wilson
lines; an individual diagram need not vanish.

\subsection{Collinear-to-hard approximation: Sec.~2.4.1}
\label{sec:notes-collinear-hard}
\subsubsection{The original connection: source Eq.~(2.55)}
Write $k_q$ and $l_{\bar q}$ for the connecting fermion momenta.
The source uses both $k_{\bar q}$ and $l_{\bar q}$ for the latter.
With indices mostly suppressed,
\begin{equation}
 C_q\otimes H\otimes C_{\bar q}
 =\int[d k][d l]\;C_q(\{k\})
 \left(\prod_i g^{\mu_i\nu_i}\right)H(\{k\};\{l\})
 \left(\prod_jg^{\sigma_j\rho_j}\right)C_{\bar q}(\{l\}),
\end{equation}
where
\begin{equation}
 \int[d k][d l]=\int\frac{d^4k_q\,d^4l_{\bar q}}{(2\pi)^8}
 \prod_{i=1}^{n}\frac{d^4k_i}{(2\pi)^4}
 \prod_{j=1}^{m}\frac{d^4l_j}{(2\pi)^4}.
\end{equation}
\textbf{$\otimes$ :both integration and index contraction.}


\subsubsection{Momentum projection: source Eqs.~(2.57)--(2.59)}
Define the lightlike vectors
\begin{align}
 n^\mu&=(1,0,0,-1)/\sqrt2,\qquad
 \bar n^\mu=(1,0,0,1)/\sqrt2,\\
 n^2&=\bar n^2=0,\qquad n\cdot\bar n=1,\qquad
 k\cdot n=k^+,\quad k\cdot\bar n=k^-.
\end{align}
Then $k^\mu=k^+\bar n^\mu+k^-n^\mu+k_T^\mu$.
The hard part is expanded in small components:
\begin{align}
 H(\{k\};\{l\})&\longmapsto H(\{\hat k\};\{\hat l\}),\\
 \hat k_i^\mu&=(k_i^+,0,\bm0_T)=(k_i\cdot n)\bar n^\mu,\\
 \hat l_j^\mu&=(0,l_j^-,\bm0_T)=(l_j\cdot\bar n)n^\mu.
\end{align}
Hard propagators are not pinched in this region. After suitable contour
deformations, they have virtualities $O(Q^2)$ and can be expanded.
For example,
\begin{equation}
 (k+l)^2=2k^+l^-+O(\lambda^2),\qquad 2k^+l^-=O(Q^2).
\end{equation}
In a collinear denominator, however, $2k^+k^-$ and $\bm k_T^2$
are both $O(\lambda^2)$ and must be kept. Thus
$C_q(k)H(k)\to C_q(k)H(\hat k)$, while $C_q$ retains the full $k$.
The projected momenta obey $\hat k^2=\hat l^2=0$.

\subsubsection{Spin projection: source Eqs.~(2.60)--(2.63)}
The leading quark spin components are selected by a projector. With
$\gamma^\pm=(\gamma^0\pm\gamma^3)/\sqrt2$,
\begin{equation}
 (\gamma^\pm)^2=0,\quad\{\gamma^+,\gamma^-\}=2,\quad
 \slsh n=\gamma^+,\quad\slsh{\hat k}=k^+\gamma^-.
\end{equation}
Using the massless spin sum,
\begin{equation}
 \mathcal P
 =\frac{\slsh n\sum_su_s(\hat k)\bar u_s(\hat k)}{2\hat k\cdot n}
 =\frac{\gamma^+k^+\gamma^-}{2k^+}
 =\boxed{\frac{\gamma^+\gamma^-}{2}}.
\end{equation}
This sum defines a projector, not a spin average. Check
\begin{align}
 \mathcal P^2&=\tfrac14\gamma^+(2-\gamma^+\gamma^-)\gamma^-
 =\mathcal P,\qquad \mathcal P\slsh{\hat k}=0,\\
 \bar u(\hat k)\mathcal P
 &=\tfrac12\bar u(\hat k)(2-\gamma^-\gamma^+)=\bar u(\hat k).
\end{align}
The last step uses $\bar u(\hat k)\gamma^-=0$.
The leading numerator $k^+\gamma^-$ is kept because
\begin{equation}
 (k^+\gamma^-)\mathcal P
 =\tfrac{k^+}{2}\gamma^-\gamma^+\gamma^-=k^+\gamma^-.
\end{equation}
On the other side,
\begin{equation}
 \overline{\mathcal P}
 =\frac{\slsh{\hat l}\slsh{\bar n}}{2\hat l\cdot\bar n}
 =\frac{l^-\gamma^+\gamma^-}{2l^-}=\mathcal P.
\end{equation}
The massless antiquark spin sum also equals $\slsh{\hat l}$.
The matrix is the same, but it acts on the other side:
\begin{equation}
 C_qHC_{\bar q}\simeq
 C_q\mathcal P H(\{\hat k\};\{\hat l\})
 \overline{\mathcal P}C_{\bar q}.
\end{equation}
With all indices and integration measures shown, source Eq.~(2.63) is
\begin{equation}
\begin{aligned}
 &C_q\otimes H\otimes C_{\bar q}\\
 &\simeq\int
 \frac{d^4k_q}{(2\pi)^4}\frac{d^4l_{\bar q}}{(2\pi)^4}
 \left[\prod_{i=1}^{n}\frac{d^4k_i}{(2\pi)^4}\right]
 \left[\prod_{j=1}^{m}\frac{d^4l_j}{(2\pi)^4}\right]\\
 &\quad\times C_{q,\alpha,\mu_1\cdots\mu_n}
 (k_q,k_1,\ldots,k_n)\,
 g^{\mu_1\nu_1}\cdots g^{\mu_n\nu_n}\\
 &\quad\times\mathcal P_{\alpha\delta}
 H_{\delta,\nu_1\cdots\nu_n;\kappa,\sigma_1\cdots\sigma_m}
 \bigl(\hat k_q,\hat k_1,\ldots,\hat k_n;\\
 &\hspace{5em}\hat l_{\bar q},\hat l_1,\ldots,\hat l_m\bigr)
 \overline{\mathcal P}_{\kappa\beta}\\
 &\quad\times g^{\sigma_1\rho_1}\cdots g^{\sigma_m\rho_m}
 C_{\bar q,\beta,\rho_1\cdots\rho_m}
 (l_{\bar q},l_1,\ldots,l_m).
\end{aligned}
\label{eq:notes-full-projected-hard}
\end{equation}
Here $\overline{\mathcal P}=\mathcal P$; we use $l_{\bar q}$
consistently for the antiquark momentum.
The connecting fermion legs of $H$ are amputated and use on-shell
spinors. Internal lines in $C$ still retain their virtualities.

\subsubsection{Gluon projection: source Eqs.~(2.64)--(2.65)}
Start with the exact identity
\begin{equation}
 g^{\mu\nu}=n^\mu\bar n^\nu+\bar n^\mu n^\nu+g_T^{\mu\nu}.
\end{equation}
For an extra $+z$ collinear gluon, power counting selects
$(C\cdot n)(\bar n\cdot H)=C^+H^-$ at leading power.
Since $\hat k^\nu=(k\cdot n)\bar n^\nu$,
\begin{equation}
 \boxed{g^{\mu\nu}\longmapsto n^\mu\bar n^\nu
 =\frac{n^\mu\hat k^\nu}{k\cdot n}.}
 \label{eq:notes-gluon-hard}
\end{equation}
The arrow is an approximation; the equality is exact. Thus
\begin{equation}
 C_\mu g^{\mu\nu}H_\nu\simeq
 \frac{C\cdot n}{k\cdot n}\,\hat k^\nu H_\nu.
\end{equation}
On the opposite side,
\begin{equation}
 g^{\sigma\rho}\longmapsto
 \frac{\hat l^\sigma\bar n^\rho}{l\cdot\bar n}.
\end{equation}
For all gluons, set $\mathcal H=\mathcal P H(\hat k,\hat l)
\overline{\mathcal P}$ and write, with indices understood,
\begin{align}
 C_q\otimes H\otimes C_{\bar q}\simeq\int[d k][d l]\;&
 C_q\prod_i\frac{n^{\mu_i}}{k_i\cdot n}
 \left(\prod_i\hat k_i^{\nu_i}\right)\mathcal H\notag\\
 &\times\left(\prod_j\hat l_j^{\sigma_j}\right)
 \prod_j\frac{\bar n^{\rho_j}}{l_j\cdot\bar n}C_{\bar q}.
\end{align}
Each gluon now contracts $H$ with its momentum. This is the form needed
for Ward identities on the hard part.

\subsection{Soft-to-collinear approximation: Sec.~2.4.2}
\label{sec:notes-soft-collinear}
\subsubsection{Expand the denominator: source Eq.~(2.66)}
A uniformly soft momentum has $k_s\sim(\lambda_S,\lambda_S,\lambda_S)$,
with $\lambda_S\ll Q$. It can be dropped inside $H$ at leading power,
but can change a small collinear denominator:
\begin{align}
 (k_i+k_s)^2-m^2&=(k_i^2-m^2)+2k_i\cdot k_s+k_s^2,\\
 2k_i\cdot k_s&=2k_i^+k_s^-+2k_i^-k_s^+
 -2\bm k_{iT}\cdot\bm k_{sT}.
\end{align}
For a $+z$ collinear momentum, the four soft-dependent terms have sizes
\begin{equation}
 Q\lambda_S,\qquad \lambda^2\lambda_S/Q,\qquad
 \lambda\lambda_S,\qquad\lambda_S^2,
\end{equation}
respectively. The first is largest, so
\begin{equation}
 \boxed{(k_i+k_s)^2-m^2+\ii0
 \simeq(k_i^2-m^2)+2k_i^+k_s^-+\ii0.}
\end{equation}
The numerator is led by $k_i^+\gamma^-=O(Q)$; the added soft term is
only $O(\lambda_S)$. Soft momentum can therefore be dropped in the
leading numerator, but must be kept in the denominator.

\subsubsection{Keep one soft component: source Eqs.~(2.67)--(2.69)}
Only $k_s^-$ is needed inside $C_q$. For $C_{\bar q}$, the leading term
is $2l^-k_s^+$, so keep $k_s^+$:
\begin{align}
 \hat k_{qs}^\mu&=(0,k_{qs}^-,\bm0_T)=(k_{qs}\cdot\bar n)n^\mu,\\
 \hat k_{\bar qs}^\mu&=(k_{\bar qs}^+,0,\bm0_T)
 =(k_{\bar qs}\cdot n)\bar n^\mu.
\end{align}
The original collinear virtuality is kept over the range
$\lambda^2/Q\lesssim\lambda_S\lesssim\lambda$.
At its lower end, $Q\lambda_S\sim\lambda^2$, so it is as important as
the soft shift. Keeping both gives one approximation across this range.

The soft connection is schematically
\begin{equation}
 \int[d k_s]\;C_q(\{\hat k_{qs}\})\,g\,
 S(\{k_{qs}\};\{k_{\bar qs}\})\,g\,
 C_{\bar q}(\{\hat k_{\bar qs}\}),
\end{equation}
where $[d k_s]$ includes $d^4k_s/(2\pi)^4$ for each connection.
$S$ includes soft propagators and momentum constraints and keeps its
full soft arguments. The full collinear arguments of $C$ are also kept,
but are not shown here.

\subsubsection{Project the gluon indices: source Eqs.~(2.70)--(2.71)}
This time place the momentum factor on the collinear end:
\begin{align}
 g^{\mu\nu}&\longmapsto
 \frac{\hat k_{qs}^\mu\bar n^\nu}{k_{qs}\cdot\bar n},\\
 g^{\sigma\rho}&\longmapsto
 \frac{n^\sigma\hat k_{\bar qs}^\rho}{n\cdot k_{\bar qs}}.
\end{align}
Both expressions equal the selected $n\bar n$ term. For example,
\begin{equation}
 C_{q,\mu}g^{\mu\nu}S_\nu\simeq
 (\hat k_{qs}^\mu C_{q,\mu})
 \frac{\bar n^\nu S_\nu}{k_{qs}\cdot\bar n}.
\end{equation}
For all soft connections, source Eq.~(2.71) reads
\begin{equation}
\begin{aligned}
 &C_q\otimes S\otimes C_{\bar q}\\
 &\simeq\int
 \left[\prod_{i=1}^{n}\frac{d^4k_{qs,i}}{(2\pi)^4}\right]
 \left[\prod_{j=1}^{m}\frac{d^4k_{\bar qs,j}}{(2\pi)^4}\right]\\
 &\quad\times C_{q;\mu_1\cdots\mu_n}
 (\{\hat k_{qs,i}\})
 \left(\prod_{i=1}^{n}\hat k_{qs,i}^{\mu_i}\right)\\
 &\quad\times\left[\prod_{i=1}^{n}
 \frac{\bar n^{\nu_i}}{k_{qs,i}\cdot\bar n}\right]
 S_{\nu_1\cdots\nu_n;\sigma_1\cdots\sigma_m}
 \bigl(\{k_{qs,i}\},\{k_{\bar qs,j}\}\bigr)\\
 &\quad\times\left[\prod_{j=1}^{m}
 \frac{n^{\sigma_j}}{n\cdot k_{\bar qs,j}}\right]
 \left(\prod_{j=1}^{m}\hat k_{\bar qs,j}^{\rho_j}\right)
 C_{\bar q;\rho_1\cdots\rho_m}(\{\hat k_{\bar qs,j}\}).
\end{aligned}
\label{eq:notes-full-soft-connections}
\end{equation}
Ward identities now act on the collinear factors. The pole prescriptions
of the later Wilson-line factors need the full derivation.

\subsection{Link to the SIDIS calculation}
The two approximations give
\begin{equation}
 \begin{array}{c|c|c}
 \text{connection}&\text{momentum contraction}&\text{Ward identity on}\\\hline
 \text{collinear to hard}&\hat k^\nu H_\nu&H\\
 \text{soft to collinear}&\hat k_s^\mu C_\mu&C
 \end{array}
\end{equation}
Sec.~2.4 prepares these identities. A full proof also needs diagram sums,
Wilson lines, and overlap subtractions. Other regions, including Glauber
momenta, need further analysis.
For SIDIS, $ea\to ea$ gives the tree-level hard part; PDFs and FFs
describe the initial and final collinear dynamics after factorization.

\clearpage
\section{Single-hadron SIDIS}
\label{sec:single-hadron}
Apply the approximations of Section~\ref{sec:factorization-notes}
to $ep\to ehX$. Keep the leading power and the tree-level hard part.
\begin{equation}
 \text{process and variables}\ \longrightarrow\
 \text{projected hard part}\ \longrightarrow\
 \text{PDF and FF}.
\end{equation}

\subsection{Process and kinematics}
\label{sec:kinematics}
Consider unpolarized single-hadron SIDIS:
\begin{equation}
 e^-(\ell)+p(P)\longrightarrow e^-(\ell')+h(P_h)+X.
 \label{eq:process}
\end{equation}
\begin{itemize}
 \item $\ell,\ell'$: incoming and scattered electron four-momenta.
 \item $P$: proton four-momentum; $P^2=M^2$.
 \item $P_h$: observed hadron four-momentum.
 \item $X$: all unobserved final particles.
\end{itemize}
Use $g^{\mu\nu}=\operatorname{diag}(1,-1,-1,-1)$ and neglect $m_e$.
Define the photon momentum and its virtuality:
\begin{equation}
 \boxed{q=\ell-\ell',\qquad Q^2=-q^2>0.}
 \label{eq:transfer}
\end{equation}
Define the measured variables:
\begin{equation}
 \boxed{x_B=\frac{Q^2}{2P\cdot q},\qquad
 y=\frac{P\cdot q}{P\cdot\ell},\qquad
 z=\frac{P\cdot P_h}{P\cdot q}.}
 \label{eq:variables}
\end{equation}
\begin{itemize}
 \item $Q$: hard scale; require $Q\gg\Lambda_{\rm QCD}$.
 \item $x_B$: Bjorken variable, fixed by electron kinematics.
 At tree level and leading power, it equals the incoming parton fraction.
 \item $y$: fraction of the electron energy transferred to the hadron system,
 in the proton rest frame.
 \item $z$: observed hadron energy divided by the transferred energy,
 in the proton rest frame.
\end{itemize}


\subsubsection{In the proton rest frame}
Choose the incoming electron along $+z$:
\begin{align}
 P&=(M,0,0,0),\qquad \ell=(E,0,0,E),\\
 \ell'&=(E',E'\sin\theta,0,E'\cos\theta),\qquad
 P_h=(E_h,\bm P_h).
\end{align}
Set $\nu=E-E'$:
\begin{equation}
 P\cdot q=M\nu,\qquad P\cdot\ell=ME,\qquad P\cdot P_h=ME_h.
\end{equation}
\begin{equation}
  \boxed{x_B=\frac{Q^2}{2M\nu},\qquad
 y=\frac{E-E'}{E},\qquad z=\frac{E_h}{E-E'}}.
\end{equation}
Use $\ell^2=\ell'^2=0$ to obtain
\begin{align}
 Q^2&=-(\ell-\ell')^2=2\ell\cdot\ell'\\
 &=2EE'(1-\cos\theta)
 =4EE'\sin^2\frac{\theta}{2}.
 \label{eq:virtuality}
\end{align}
\begin{figure}[H]
 \centering
 \includegraphics[width=0.70\textwidth]{figures/sidis_momentum_triangle.png}
 \caption{Three-momentum relation: $\bm q=\bm\ell-\bm\ell'$.}
 \label{fig:triangle}
\end{figure}

\subsection{The projected tree-level hard part}
\label{sec:hard-scattering}
\textbf{Identify the subgraphs as in Section~\ref{sec:factorization-foundations}.}
\begin{equation}
 \underbrace{C_{\rm in}}_{\text{proton direction}}
 \longleftrightarrow H\longleftrightarrow
 \underbrace{C_{\rm out}}_{\text{observed-hadron direction}}.
\end{equation}
\begin{itemize}
 \item $C_{\rm in}$: the initial collinear subgraph in the proton.
 \item $H$: the photon vertex; include the electron current in its hard coefficient.
 \item $C_{\rm out}$: the final collinear subgraph containing $h$ and other particles.
 \item Assign the connecting collinear propagators to $C$, as in Section 1.
\end{itemize}

\subsubsection{Momentum and spin projections}
Use a fast frame with the proton along $+z$ and the current jet along $-z$.
Use the same $n,\bar n$ as in Sec.~\ref{sec:notes-collinear-hard}:
\begin{align}
 k&\sim(Q,\lambda^2/Q,\lambda),&
 r&\sim(\lambda^2/Q,Q,\lambda),\\
 \hat k&=(k^+,0,\bm0_T),&
 \hat r&=(0,r^-,\bm0_T).
\end{align}
Apply the approximation \textbf{only to the hard part}:
\begin{equation}
 C_{\rm in}(k)H(k,r)C_{\rm out}(r)
 \longmapsto
 C_{\rm in}(k)H(\hat k,\hat r)C_{\rm out}(r).
\end{equation}
This equation shows the momentum dependence only; spin and gluon
connections are treated as in Section 1.
\begin{align}
 \mathcal P&=\frac{\gamma^+\gamma^-}{2},&
 \sum_su_s(\hat k)\bar u_s(\hat k)&=\slsh{\hat k},\\
 \hat k^2&=\hat r^2=0,&
 \slsh{\hat k}&=k^+\gamma^-,\quad
 \slsh{\hat r}=r^-\gamma^+.
\end{align}
\begin{itemize}
 \item Use $\mathcal P$ on the appropriate hard-facing Dirac index,
 as in Section 1. Extract the on-shell spinors in $H$.
 \item \textbf{Do not confuse the leading-component projector with a spin average.}
 Average the external spins only when forming the unpolarized hard coefficient.
 \item In the calculation below, write $k=\hat k$, $k'=\hat r$ for brevity.
 Thus $k^2=k'^2=0$ applies to the hard part, not to every line in the proton.
\end{itemize}
The hard process is
\begin{equation}
 e^-(\ell)+a(k)\longrightarrow e^-(\ell')+a(k'),\qquad
 q=\ell-\ell'=k'-k.
\end{equation}
\begin{itemize}
 \item $a$: a quark flavor; \textbf{$e_u=2/3$, $e_d=-1/3$.}
 \item Neglect $m_e$ and the quark mass in the hard part.
 \item Sum over final spins and color. Average over initial spins and quark color.
\end{itemize}


\subsubsection{Partonic invariants}
Use momentum conservation and the massless on-shell conditions:
\begin{align}
 \ell+k&=\ell'+k',\qquad
 \ell^2=\ell'^2=k^2=k'^2=0,\\
 \hat s&=(\ell+k)^2=2\ell\cdot k=2\ell'\cdot k',\\
 \hat t&=(\ell-\ell')^2=-2\ell\cdot\ell'=-2k\cdot k'=-Q^2,\\
 \hat u&=(\ell-k')^2=-2\ell\cdot k'=-2\ell'\cdot k.
 \label{eq:hard-dot-products}
\end{align}
\begin{equation}
 \hat s+\hat t+\hat u
 =2\ell\cdot(k-\ell'-k')=-2\ell^2=0.
 \label{eq:hard-mandelstam}
\end{equation}

\subsubsection{The amplitude and current conservation}
Take $e>0$. Use charges $-e$ and $ee_a$:
\begin{align}
 \ii\M={}&
 [\bar u(\ell')(+\ii e\gamma^\mu)u(\ell)]
 \frac{-\ii g_{\mu\nu}}{q^2}
 [\bar u(k')(-\ii ee_a\gamma^\nu)u(k)]\\
 ={}&-\frac{\ii e^2e_a}{q^2}
 [\bar u(\ell')\gamma^\mu u(\ell)]
 [\bar u(k')\gamma_\mu u(k)].
\end{align}
Remove the overall $\ii$ and use $q^2=-Q^2$:
\begin{equation}
 \boxed{\M=\frac{e^2e_a}{Q^2}
 [\bar u(\ell')\gamma^\mu u(\ell)]
 [\bar u(k')\gamma_\mu u(k)].}
 \label{eq:hard-amplitude}
\end{equation}
Keep the color factor separate: $\M_{c'c}=\delta_{c'c}\M$.

Check the tree-level identity from Sec.~\ref{sec:notes-ward-identity}:
\begin{align}
 q_\mu J_e^\mu
 &=\bar u(\ell')(\slsh\ell-\slsh{\ell'})u(\ell)=0,\\
 q_\mu J_a^\mu
 &=\bar u(k')(\slsh{k'}-\slsh k)u(k)=0.
\end{align}
Use the massless Dirac equations. Both currents are conserved.

\subsubsection{Spin sums and color}
Use
\begin{equation}
 \sum_su(p,s)\bar u(p,s)=\slsh p.
\end{equation}
For example, sum the electron spins:
\begin{align}
 &\sum_{s,s'}[\bar u(\ell',s')\gamma^\mu u(\ell,s)]
 [\bar u(\ell,s)\gamma^\nu u(\ell',s')]\notag\\
 &\qquad=\operatorname{Tr}
 [\slsh{\ell'}\gamma^\mu\slsh\ell\gamma^\nu]
 \equiv L_e^{\mu\nu}.
 \label{eq:hard-current-trace}
\end{align}
Do the same for the quark:
\begin{equation}
 L^{(a)}_{\mu\nu}
 =\operatorname{Tr}[\slsh{k'}\gamma_\mu\slsh k\gamma_\nu].
\end{equation}
Include the averages:
\begin{equation}
 \underbrace{\frac12\frac12}_{\text{initial spins}}=\frac14,
 \qquad
 \underbrace{\frac1{N_c}\sum_{c,c'}|\delta_{c'c}|^2}_{\text{color}}=1.
\end{equation}
Hence
\begin{equation}
 \boxed{\overline{|\M|^2}
 =\frac{e^4e_a^2}{4Q^4}L_e^{\mu\nu}L^{(a)}_{\mu\nu}.}
 \label{eq:hard-trace-start}
\end{equation}
\textbf{The photon preserves the active quark color; $N_c=3$ in QCD.}

\subsubsection{Trace contraction}
Apply
\begin{equation}
 \operatorname{Tr}(\gamma^\alpha\gamma^\mu\gamma^\beta\gamma^\nu)
 =4(g^{\alpha\mu}g^{\beta\nu}-g^{\alpha\beta}g^{\mu\nu}
 +g^{\alpha\nu}g^{\mu\beta}).
\end{equation}
Obtain
\begin{align}
 L_e^{\mu\nu}
 &=4[\ell'^\mu\ell^\nu+\ell'^\nu\ell^\mu
 -g^{\mu\nu}(\ell'\cdot\ell)],
 \label{eq:hard-electron-trace}\\
 L^{(a)}_{\mu\nu}
 &=4[k'_\mu k_\nu+k'_\nu k_\mu-g_{\mu\nu}(k'\cdot k)].
 \label{eq:hard-quark-trace}
\end{align}
Set $c=\ell'\cdot\ell$ and $d=k'\cdot k$. Expand the product:
\begin{align}
 \frac{L_e^{\mu\nu}L^{(a)}_{\mu\nu}}{16}
 ={}& (\ell'\cdot k')(\ell\cdot k)
 +(\ell'\cdot k)(\ell\cdot k')\notag\\
 &+(\ell'\cdot k)(\ell\cdot k')
 +(\ell'\cdot k')(\ell\cdot k)\notag\\
 &-2cd-2cd+4cd\\
 ={}&2[(\ell'\cdot k')(\ell\cdot k)
 +(\ell'\cdot k)(\ell\cdot k')].
 \label{eq:hard-tensor-contraction}
\end{align}
The metric terms cancel because $g^{\mu\nu}g_{\mu\nu}=4$.
Insert the result:
\begin{align}
 \overline{|\M|^2}
 &=\frac{8e^4e_a^2}{Q^4}
 [(\ell'\cdot k')(\ell\cdot k)
 +(\ell'\cdot k)(\ell\cdot k')]
 \label{eq:hard-dot-result}\\
 &=\frac{8e^4e_a^2}{Q^4}
 \left(\frac{\hat s^2}{4}+\frac{\hat u^2}{4}\right)\\
 &=\boxed{2e^4e_a^2\frac{\hat s^2+\hat u^2}{\hat t^2}}.
 \label{eq:hard-mandelstam-result}
\end{align}

\subsubsection{Flux and two-body phase space}
Use the massless incident flux $4\ell\cdot k=2\hat s$:
\begin{equation}
 d\hat\sigma=\frac{\overline{|\M|^2}}{2\hat s}\,d\Phi_2,
\end{equation}
where
\begin{equation}
 d\Phi_2=(2\pi)^4\delta^{(4)}(\ell+k-\ell'-k')
 \frac{d^3\bm\ell'}{(2\pi)^3\,2E_{\ell'}}
 \frac{d^3\bm k'}{(2\pi)^3\,2E_{k'}}.
 \label{eq:hard-phase-space}
\end{equation}
In the partonic centre-of-momentum frame, set $r=|\bm\ell'|=|\bm k'|$.
Integrate the delta functions:
\begin{align}
 d\Phi_2&=\frac{d\Omega}{16\pi^2}
 \int_0^\infty dr\,\delta(\sqrt{\hat s}-2r)
 =\frac{d\Omega}{32\pi^2},\\
 \frac{d\hat\sigma}{d\Omega}
 &=\frac{\overline{|\M|^2}}{64\pi^2\hat s},\qquad
 \frac{d\hat\sigma}{d\cos\hat\theta}
 =\frac{\overline{|\M|^2}}{32\pi\hat s}.
\end{align}
Here $\hat\theta$ is the electron angle in this partonic frame.
It differs from the proton-rest-frame angle $\theta$.
Use
\begin{equation}
 \hat t=-\frac{\hat s}{2}(1-\cos\hat\theta),\qquad
 \frac{d\hat t}{d\cos\hat\theta}=\frac{\hat s}{2}.
 \label{eq:hard-cm-invariants}
\end{equation}
Change variables and insert the squared amplitude:
\begin{align}
 \frac{d\hat\sigma}{d\hat t}
 &=\frac{\overline{|\M|^2}}{32\pi\hat s}\frac{2}{\hat s}\\
 &=\boxed{\frac{2\pi\alpha^2e_a^2}{\hat s^2}
 \frac{\hat s^2+\hat u^2}{\hat t^2}},\qquad
 \alpha=\frac{e^2}{4\pi}.
 \label{eq:hard-dsigma-dt}
\end{align}

\subsubsection{The hard coefficient in \texorpdfstring{$y$}{y}}
For a collinear incoming parton, use
\begin{equation}
 y=\frac{P\cdot q}{P\cdot\ell}
 \simeq\frac{k\cdot q}{k\cdot\ell}.
\end{equation}
Apply $k'=k+q$ and $k'^2=k^2=0$:
\begin{align}
 0&=(k+q)^2=2k\cdot q-Q^2,\\
 y&=\frac{Q^2}{\hat s}=-\frac{\hat t}{\hat s},\qquad
 \hat u=-\hat s-\hat t=-\hat s(1-y).
 \label{eq:hard-y-relations}
\end{align}
Then
\begin{equation}
 \frac{\hat s^2+\hat u^2}{\hat t^2}
 =\frac{1+(1-y)^2}{y^2},\qquad
 \left|\frac{d\hat t}{dy}\right|_{\hat s}=\hat s.
\end{equation}
Obtain the hard cross section:
\begin{align}
 \frac{d\hat\sigma_a}{dy}
 &=\frac{d\hat\sigma_a}{d\hat t}\,\hat s\\
 &=\frac{2\pi\alpha^2e_a^2}{\hat s y^2}[1+(1-y)^2]\\
 &=\boxed{\frac{2\pi\alpha^2e_a^2}{Q^2y}[1+(1-y)^2]}.
 \label{eq:hard-dsigma-dy}
\end{align}

\subsubsection{Color factors}
Table~\ref{tab:color} collects the color quantum numbers of the four
particles that enter this calculation.

\begin{table}[H]
\centering
\vspace*{2mm}
\caption{Color of the particles that enter this calculation. Here $N_c=3$ is
the dimension of the fundamental representation of $\mathrm{SU}(3)$; it is a
parameter of the theory, not a property of a single particle. What enters an
average over an initial state is the dimension $d_R$ of the representation
that the particle carries.}
\label{tab:color}
\small
\begin{tabular}{@{}lcccc@{}}
\toprule
particle & color representation & dimension $d_R$
& average $1/d_R$ & color factor \\
\midrule
photon $\gamma$ & singlet $\mathbf 1$
  & $1$ & --- & $1$ \\
proton $p$ & singlet $\mathbf 1$
  & $1$ & --- & $1$ \\
quark $q$ & fundamental $\mathbf 3$
  & $N_c=3$ & $1/N_c=1/3$ & $1$ \\
gluon $g$ & adjoint $\mathbf 8$
  & $N_c^2-1=8$ & $1/(N_c^2-1)=1/8$ & $C_F=4/3$ \\
\bottomrule
\end{tabular}
\end{table}

\textbf{How to read Table~\ref{tab:color}.}
\begin{itemize}
 \item A photon is a color singlet, so its vertex carries no $T^a$ and color
       space is trivial, $\M_{c'c}=\delta_{c'c}\M$. This is the $\delta_{c'c}$
       of Eq.~\eqref{eq:hard-amplitude}, and it gives the color factor $1$.
 \item A proton is a color singlet only after its constituents are combined,
       $|p\rangle\propto\varepsilon_{ijk}q^iq^jq^k$. Because it is a singlet,
       the color of the active quark is uniformly distributed,
       $\langle\bar q_iq_j\rangle\propto\delta_{ij}/N_c$, which is what makes
       the average $1/N_c$ legitimate. The two singlet rows are not the same
       statement: the photon carries no color, the proton carries balanced
       color.
 \item The entry $1$ is for singlet exchange. If a gluon is exchanged, the
       vertex is $(T^a)_{c'c}$ and the color factor changes:
       \begin{equation}
        \frac{1}{N_c}\sum_{a}\operatorname{tr}(T^aT^a)
        =\frac{T_F}{N_c}\left(N_c^2-1\right)=C_F=\frac43,
        \label{eq:color-casimir}
       \end{equation}
       with $T_F=\tfrac12$ and $C_A=N_c=3$. Note that $C_A$ and
       $d_{\text{fund}}$ both coincide numerically with $N_c$ but are
       different objects. An initial gluon is averaged with
       $1/(N_c^2-1)=1/8$; for $\gamma^*g\to q\bar q$ the color factor is
       $T_F=1/2$.
\end{itemize}

\subsection{The factorized SIDIS cross section}
\label{sec:sidis-factorization}
\subsubsection{From subgraphs to PDFs and FFs}

Use the amplitude \textbf{and its conjugate}, and sum over $X$:
\begin{equation}
 W^{\mu\nu}\ \propto\ \sum_X
 \langle P|j^\mu|hX\rangle\langle hX|j^\nu|P\rangle.
\end{equation}
The sum includes the final-state phase space and momentum constraint.
The two collinear factors now become cut matrix elements:
\begin{equation}
 \underbrace{\Phi}_{\text{proton matrix element}}
 \ \otimes\ \underbrace{\mathcal H}_{\text{hard amplitude squared}}
 \ \otimes\ \underbrace{\Delta}_{\text{fragmentation matrix element}}.
\end{equation}
\begin{itemize}
 \item Project the leading unpolarized Dirac components of $\Phi,\Delta$.
 Their scalar functions are $f_1,D_1$.
 \item Keep their low-virtuality dynamics inside these functions.
 \item Use the calculated $d\hat\sigma_a/dy$ for the hard coefficient.
\end{itemize}

\begin{figure}[H]
 \centering
 \includegraphics[width=0.9\textwidth]{figures/sidis_factorization.png}
 \caption{PDF $\otimes$ hard scattering $\otimes$ FF.}
 \label{fig:sidis}
\end{figure}
\begin{itemize}
 \item \textbf{PDF:} use $f_1^a(\xi,\mu)$ for a parton in the proton.
 Define the incoming fraction $\xi=k^+/P^+$.
 \item \textbf{Hard part:} use the calculated $d\hat\sigma_a/dy$.
 \item \textbf{FF:} use $D_1^{a\to h}(\zeta,\mu)$ for a hadron from
 the outgoing parton. In a frame with the jet near $-z$, define
 $\zeta=P_h^-/k'^-$.
 \item Choose the factorization scale $\mu\sim Q$.
 Treat the PDF and FF as nonperturbative inputs.
\end{itemize}
\textbf{Use the approximations from Section~\ref{sec:factorization-notes}:}
\begin{enumerate}
 \item Drop the small collinear components \textbf{only in the hard part}.
 Use on-shell projected partons; retain low virtualities in the collinear factors.
 \item Apply Ward identities to the extra longitudinal gluons.
 Collect their couplings into the gauge links of PDFs and FFs.
 \item Use the integrated SIDIS factorization theorem for the soft terms.
 After the required cancellations and overlap treatment, no separate
 measured soft function appears in this collinear formula.
\end{enumerate}
Section~2.4 supplies the local approximations. The full soft cancellation
and overlap treatment are additional parts of the SIDIS factorization theorem.


\subsubsection{Fraction matching}
\textbf{Incoming fraction.} Use the lightlike projection
\begin{equation}
 \hat k^\mu=\xi P^+\bar n^\mu,\qquad\hat k^2=0.
 \label{eq:sidis-hard-projections}
\end{equation}
Write $\hat k\simeq\xi P$ in leading-power scalar products.
Apply the outgoing on-shell condition:
\begin{align}
 0&=(\hat k+q)^2\simeq2\xi P\cdot q-Q^2\\
 &=Q^2\left(\frac{\xi}{x_B}-1\right)
 \quad\Rightarrow\quad \boxed{\xi=x_B}.
\end{align}
\textbf{Outgoing fraction.} Use $P_h\simeq\zeta k'$ at leading power:
\begin{align}
 z&=\frac{P\cdot P_h}{P\cdot q}
 \simeq\zeta\frac{P\cdot k'}{P\cdot q}\\
 &=\zeta\left(1+\frac{P\cdot\hat k}{P\cdot q}\right)
 \simeq\zeta
 \quad\Rightarrow\quad \boxed{\zeta=z}.
\end{align}
The omitted term gives a target-mass correction. Keep physical hadron
masses; drop their corrections only in the hard kinematics.
\begin{equation}
 \boxed{\xi=x_B,\qquad\zeta=z\quad
 \text{at tree level and leading power}.}
 \label{eq:sidis-fraction-matching}
\end{equation}

\subsubsection{The tree-level convolution}
Write the tree-level number-density form:
\begin{align}
 \frac{d\sigma}{dx_B\,dy\,dz}
 =\sum_{a=q,\bar q}\int_0^1d\xi\int_0^1d\zeta\;&
 f_1^a(\xi,\mu)\frac{d\hat\sigma_a(\xi)}{dy}
 D_1^{a\to h}(\zeta,\mu)\notag\\
 &\times\delta(x_B-\xi)\delta(z-\zeta).
 \label{eq:sidis-lo-convolution}
\end{align}
Use the delta functions:
\begin{align}
 \int_0^1d\xi\,f_1^a(\xi,\mu)
 \frac{d\hat\sigma_a(\xi)}{dy}\delta(x_B-\xi)
 &=f_1^a(x_B,\mu)
 \left.\frac{d\hat\sigma_a}{dy}\right|_{\xi=x_B},\\
 \int_0^1d\zeta\,D_1^{a\to h}(\zeta,\mu)\delta(z-\zeta)
 &=D_1^{a\to h}(z,\mu).
\end{align}
Therefore
\begin{equation}
 \boxed{\frac{d\sigma}{dx_B\,dy\,dz}
 =\sum_{a=q,\bar q}f_1^a(x_B,\mu)
 \left.\frac{d\hat\sigma_a}{dy}\right|_{\xi=x_B}
 D_1^{a\to h}(z,\mu).}
 \label{eq:sidis-product}
\end{equation}
\begin{itemize}
 \item Sum over quarks and antiquarks. Their unpolarized hard factors
 have the same form, weighted by $e_a^2$.
 \item \textbf{Do not repeat the spin or color average.}
 It is already in $d\hat\sigma_a/dy$.
 \item The delta functions are specific to this tree-level kinematics.
 Higher-order radiation gives nontrivial convolutions.
\end{itemize}

\subsubsection{The SIDIS cross section}
Use $s=(\ell+P)^2=M^2+2P\cdot\ell$ and check
\begin{align}
 \hat s&\simeq\xi(s-M^2),\\
 Q^2&=2x_BP\cdot q=x_By(s-M^2).
 \label{eq:sidis-Q2-xy}
\end{align}
At $\xi=x_B$, $\hat s y=Q^2$. Insert
Eq.~\eqref{eq:hard-dsigma-dy} into Eq.~\eqref{eq:sidis-product}:
\begin{equation}
 \boxed{\begin{aligned}
 \frac{d\sigma}{dx_B\,dy\,dz}
 ={}&\frac{2\pi\alpha^2}{Q^2y}[1+(1-y)^2]\\
 &\times\sum_{a=q,\bar q}e_a^2
 f_1^a(x_B,\mu)D_1^{a\to h}(z,\mu).
 \end{aligned}}
 \label{eq:sidis-final-lo}
\end{equation}


\subsubsection{Structure functions}
Use $Y_+=1+(1-y)^2$ and write
\begin{equation}
 \frac{d\sigma}{dx_B\,dy\,dz}
 =\frac{2\pi\alpha^2}{x_ByQ^2}
 [Y_+F_{UU,T}+2(1-y)F_{UU,L}].
 \label{eq:sidis-structure-cross-section}
\end{equation}
Compare with Eq.~\eqref{eq:sidis-final-lo}:
\begin{align}
 F_{UU,T}&=x_B\sum_{a=q,\bar q}e_a^2
 f_1^a(x_B,\mu)D_1^{a\to h}(z,\mu),\\
 F_{UU,L}&=0\qquad\text{at tree level}.
 \label{eq:sidis-structure-lo}
\end{align}


\subsubsection{Change of measured variable}
At fixed $x_B$ and beam energy, use
\begin{equation}
 \left|\frac{dy}{dQ^2}\right|
 =\frac1{x_B(s-M^2)}=\frac{y}{Q^2}.
\end{equation}
Multiply Eq.~\eqref{eq:sidis-final-lo} by this Jacobian:
\begin{equation}
 \boxed{\begin{aligned}
 \frac{d\sigma}{dx_B\,dQ^2\,dz}
 ={}&\frac{2\pi\alpha^2}{Q^4}[1+(1-y)^2]\\
 &\times\sum_{a=q,\bar q}e_a^2
 f_1^a(x_B,\mu)D_1^{a\to h}(z,\mu).
 \end{aligned}}
 \label{eq:sidis-Q2-cross-section}
\end{equation}

\clearpage
\section{Transverse-spin dihadron SIDIS}
\label{sec:dihadron}
Replace the single-hadron final factor by a pair factor, and keep the
transverse-spin components of the collinear matrix elements.
\begin{equation}
 C_{\rm in}(P,S_T)\ \longleftrightarrow\ H
 \ \longleftrightarrow\ C_{\rm out}(P_h,R).
\end{equation}
\begin{itemize}
 \item $C_{\rm in}$: collinear part of the polarized proton.
 \item $H$: short-distance electron--quark scattering.
 \item $C_{\rm out}$: collinear part that forms the hadron pair.
\end{itemize}
The momentum approximation is unchanged. The spin projection changes.

\subsection{Process and pair kinematics}
Consider an unpolarized electron and a transversely polarized proton:
\begin{equation}
 e^-(\ell)+p(P,S_T)\longrightarrow
 e^-(\ell')+h_1(P_1)+h_2(P_2)+X.
\end{equation}
\begin{figure}[H]
 \centering
 \includegraphics[width=0.98\textwidth]{figures/dihadron_sidis_selected.png}
 \caption{Transverse-spin dihadron SIDIS. }
 \label{fig:dihadron-sidis}
\end{figure}
Set $q=\ell-\ell'$, $Q^2=-q^2$, 
\begin{equation}
x_B=\frac{Q^2}{2P\cdot q}, \qquad y=\frac{P \cdot q}{P \cdot \ell}
\end{equation}
Define
\begin{align}
 P_h&=P_1+P_2,\qquad R=\frac{P_1-P_2}{2},\qquad M_h^2=P_h^2,\\
 z_i&=\frac{P\cdot P_i}{P\cdot q},\qquad
 z=z_1+z_2=\frac{P\cdot P_h}{P\cdot q}.
\end{align}
\begin{itemize}
 \item $P_h$: total pair momentum; $M_h$: pair invariant mass.
 \item $R$: relative momentum. 
 \item $z_i$: transferred-energy fraction carried by hadron $i$.
\end{itemize}



\subsubsection{Azimuthal angles}
For a $\pi^+\pi^-$ pair, set $P_1=P_{\pi^+}$ and $P_2=P_{\pi^-}$. 
\begin{itemize}
  \item $\vartheta$: the polar angle of $\pi^+$ in the pair rest frame;$c=\cos\vartheta$. 
  \item $\phi_R$: the azimuthal angle of the hadrons relative momentum. 
  \item $\phi_S$: the azimuthal angle of the proton transverse spin.
\end{itemize}
At leading power,
\begin{equation}
 \begin{aligned}
 |\bm R_T|&=|\bm R^*|\sqrt{1-c^2}=|\bm R^*|\sin\vartheta,\\
 |\bm R^*|&=
 \frac{\sqrt{[M_h^2-(m_1+m_2)^2][M_h^2-(m_1-m_2)^2]}}{2M_h}.
 \end{aligned}
\end{equation}
\begin{figure}[H]
 \centering
 \includegraphics[width=0.8\textwidth]{figures/dihadron_angles.png}
 \caption{Dihadron SIDIS angles. In this figure, $k,k'$ denote the
 electron momenta $\ell,\ell'$ of our notation; the inset angle $\theta$
 is our pair polar angle $\vartheta$.}
 \label{fig:dihadron-angles}
\end{figure}
\begin{itemize}
 \item View the transverse plane with $+z$ along $\bm q$ and the lepton plane as the reference plane.
 \item Set $\Phi=\phi_R+\phi_S$.
\end{itemize}


\subsection{The projected tree-level hard part with transverse spins}

Let $\lambda\ll Q$ be a small hadronic momentum scale.

Incoming quark:
\begin{align}
 k&=(k^+,k^-,\bm k_T)
   \sim(Q,\lambda^2/Q,\lambda),&
 k^2&=2k^+k^- -\bm k_T^2=O(\lambda^2).
\end{align}
Here $k^+\sim Q$ is large; $k^-$ and $|\bm k_T|$ are small compared
with $Q$.

Outgoing quark:
\begin{align}
k'&=(k'^+,k'^-,\bm k'_T)
   \sim(\lambda^2/Q,Q,\lambda),&
k'^2&=2k'^+k'^- -\bm k_T'^2=O(\lambda^2).
\end{align}
It moves along the opposite direction, so $k'^-\sim Q$ is large.
\textbf{Keep only the large components in $H$:}
\begin{align}
 \hat k&=k^+\bar n=(k^+,0,\bm0_T),\\
 \hat k'&=k'^-n=(0,k'^-,\bm0_T),\\
 \hat k^2&=\hat k'^2=0,\\
 C_{\rm in}(k)H(k,k')&C_{\rm out}(k')
\longmapsto C_{\rm in}(k)H(\hat k,\hat k')C_{\rm out}(k').
\end{align}

The hard process is
\begin{equation}
 e^-(\ell)+a(k)\longrightarrow e^-(\ell')+a(k'),\qquad
 q=\ell-\ell'=k'-k.
\end{equation}
Use the same amplitude:
\begin{equation}
  i\M=[\bar u(\ell')\gamma^\mu u(\ell)]-i\frac{e^2e_a}{q^2}[\bar u(k')\gamma_\mu u(k)].
\end{equation}
\begin{equation}
 \M=\frac{e^2e_a}{Q^2}
 [\bar u(\ell')\gamma^\mu u(\ell)]
 [\bar u(k')\gamma_\mu u(k)].
\end{equation}
\begin{align}
 \hat s&=(\ell+k)^2=2\ell\cdot k,\\
 \hat t&=(\ell-\ell')^2=q^2=-Q^2,\\
 \hat u&=(\ell-k')^2=-2\ell\cdot k'.
\end{align}
For the projected massless hard momenta,
$\hat s+\hat t+\hat u=0$. These variables will enter the spin trace.


\subsubsection{Keep the transverse-spin density matrices}


Use unit spacelike spin vectors $s,s'$ with
\begin{equation}
 s^2=s'^2=-1,\quad s\cdot k=0,\quad s'\cdot k'=0.
\end{equation}
Replace the unpolarized density $\rho_U(k)=\slsh k/2$ by
the massless limit of the transverse-spin density:
\begin{equation}
 u(k,s)\bar u(k,s)=\frac12\slsh k(1+\gamma_5\slsh s).
\end{equation}
Average only the initial electron spin. Sum only the final electron spin:
\begin{equation}
 \overline{|\M|^2}(s,s')
 =\frac{e^4e_a^2}{8Q^4}L_e^{\mu\nu}
 \operatorname{Tr}\!\left[
 \slsh{k'}(1+\gamma_5\slsh{s'})\gamma_\mu
 \slsh k(1+\gamma_5\slsh s)\gamma_\nu\right].
\end{equation}
\begin{itemize}
 \item $1/2$: electron spin average.
 \item $1/2\times1/2$: the two fixed quark-spin projectors.
 \item Color average and sum: one, as in Section 2.
\end{itemize}

\subsubsection{Calculate the squared amplitude}

Choose the partonic scattering plane as the $xz$-plane. Its normal is
\begin{equation}
 N=(0,0,1,0),\qquad s=\eta N,\quad s'=\eta'N,
 \qquad \eta,\eta'=\pm1.
\end{equation}
All external momenta lie in the scattering plane:
$N\cdot\ell=N\cdot\ell'=N\cdot k=N\cdot k'=0$.


Use $Q^4=\hat t^2$ and $\{\gamma_5,\gamma^\mu\}=0$:
\begin{align}
 \overline{|\M|^2}(\eta,\eta')
 =\frac{e^4e_a^2}{8\hat t^2}L_e^{\mu\nu}
 \Big[
 \operatorname{Tr}(\slsh{k'}\gamma_\mu\slsh k\gamma_\nu)
 \notag
 +\eta\eta'\operatorname{Tr}
 (\slsh{k'}\gamma_5\slsh N\gamma_\mu
  \slsh k\gamma_5\slsh N\gamma_\nu)\Big].
\end{align}
\textbf{The first trace: }
\begin{align}
 \operatorname{Tr}(\slsh{k'}\gamma_\mu\slsh k\gamma_\nu)
 &=4\bigl(k'_\mu k_\nu+k'_\nu k_\mu
          -g_{\mu\nu}\,k'\cdot k\bigr),\\
 L_e^{\mu\nu}
 \operatorname{Tr}(\slsh{k'}\gamma_\mu\slsh k\gamma_\nu)
 &=32\bigl[(\ell'\cdot k')(\ell\cdot k)
          +(\ell'\cdot k)(\ell\cdot k')\bigr]\notag\\
 &=8(\hat s^2+\hat u^2).
\end{align}
\textbf{The second trace:}

Move the two $\gamma_5$ matrices together. Since $\gamma_5^2=1$, $\slsh N^2=-1$, $N\cdot k=N\cdot k'=0$:
\begin{align}
 \operatorname{Tr}(\slsh{k'}\gamma_5\slsh N\gamma_\mu
                   \slsh k\gamma_5\slsh N\gamma_\nu)
 &=-\operatorname{Tr}(\slsh{k'}\slsh N\gamma_\mu
                      \slsh k\slsh N\gamma_\nu),\\
 \slsh N\gamma_\mu\slsh k\slsh N
 &=2N_\mu\slsh k\slsh N-\gamma_\mu\slsh k,\\
 \operatorname{Tr}(\slsh{k'}\slsh k\slsh N\gamma_\nu)
 &=4(k'\cdot k)N_\nu.
\end{align}
Contract with $L_e^{\mu\nu}$. Use
$L_e^{\mu\nu}N_\mu N_\nu=-2\hat t$ and
$k'\cdot k=-\hat t/2$:
\begin{align}
 &L_e^{\mu\nu}
 \operatorname{Tr}(\slsh{k'}\gamma_5\slsh N\gamma_\mu
                   \slsh k\gamma_5\slsh N\gamma_\nu)\notag\\
 &\quad=8(\hat s^2+\hat u^2)
       -8(k'\cdot k)L_e^{\mu\nu}N_\mu N_\nu\notag\\
 &\quad=8(\hat s^2+\hat u^2)-8\hat t^2\notag\\
 &\quad=-16\hat s\hat u,
 \qquad(\hat s+\hat t+\hat u=0).
\end{align}
\begin{align}
 \overline{|\M|^2}(\eta,\eta')
 &=\frac{e^4e_a^2}{8\hat t^2}
   \bigl[8(\hat s^2+\hat u^2)
         -16\eta\eta'\hat s\hat u\bigr]\notag\\
 &=\boxed{\frac{e^4e_a^2}{\hat t^2}
   \bigl[\hat s^2+\hat u^2-2\eta\eta'\hat s\hat u\bigr]}.
\end{align}
Check by summing over the final quark spin:
\begin{equation}
 \sum_{\eta'=\pm1}\overline{|\M|^2}(\eta,\eta')
 =\frac{2e^4e_a^2}{\hat t^2}(\hat s^2+\hat u^2).
\end{equation}
\subsubsection{The spin-transfer coefficient}
Fix the incoming quark spin at $\eta=+1$.

\textbf{Spin preserved:} set $\eta'=+1$, so $\eta\eta'=+1$:
\begin{equation}
 \overline{|\M|^2}(+,+)
 =\frac{e^4e_a^2}{\hat t^2}
  (\hat s^2+\hat u^2-2\hat s\hat u).
\end{equation}
\textbf{Spin reversed:} set $\eta'=-1$, so $\eta\eta'=-1$:
\begin{equation}
 \overline{|\M|^2}(+,-)
 =\frac{e^4e_a^2}{\hat t^2}
  (\hat s^2+\hat u^2+2\hat s\hat u).
\end{equation}
The two rates have the same flux and two-body phase space.
These factors cancel in their ratio. Subtract and add the squared
amplitudes:
\begin{align}
 \overline{|\M|^2}(+,+)-\overline{|\M|^2}(+,-)
 &=\frac{e^4e_a^2}{\hat t^2}(-4\hat s\hat u),\\
 \overline{|\M|^2}(+,+)+\overline{|\M|^2}(+,-)
 &=\frac{e^4e_a^2}{\hat t^2}
   2(\hat s^2+\hat u^2).
\end{align}
Therefore
\begin{equation}
 D_{NN}
 =\frac{d\hat\sigma(+,+)-d\hat\sigma(+,-)}
        {d\hat\sigma(+,+)+d\hat\sigma(+,-)}
 =\frac{-4\hat s\hat u}{2(\hat s^2+\hat u^2)}
 =\boxed{\frac{-2\hat s\hat u}{\hat s^2+\hat u^2}}.
\end{equation}
\begin{equation}
-1\leq D_{NN}\leq 1.
\end{equation}
\begin{itemize}
    \item $D_{NN}=1$: complete spin preservation.
    \item $D_{NN}=0$: no net preference for either spin direction.
    \item $D_{NN}=-1$: complete spin reversal.
\end{itemize}
Use $\hat u=-\hat s(1-y)$:
\begin{equation}
 \boxed{D_{NN}(y)=\frac{2(1-y)}{1+(1-y)^2}.}
 \label{eq:dihadron-DNN}
\end{equation}
\begin{itemize}
 \item At $y\to0$, $D_{NN}\to1$; at $y\to1$, $D_{NN}\to0$.
 \item The unpolarized hard weight is $Y_+=1+(1-y)^2$.
 The transverse-transfer weight is $Y_+D_{NN}=2(1-y)$.
\end{itemize}
Define the normal-spin correlation as the difference of the two rates:
\begin{align}
 \frac{d\Delta_T\hat\sigma_a}{dy}
 &\equiv\frac{d\hat\sigma_a(+,+)}{dy}
       -\frac{d\hat\sigma_a(+,-)}{dy},\\
 \frac{d\hat\sigma_a}{dy}
 &=\frac{d\hat\sigma_a(+,+)}{dy}
  +\frac{d\hat\sigma_a(+,-)}{dy}.
\end{align}
\begin{align}
 \frac{d\Delta_T\hat\sigma_a}{dy}
 &=D_{NN}(y)\frac{d\hat\sigma_a}{dy}\notag\\
 &=\frac{2(1-y)}{1+(1-y)^2}
   \frac{2\pi\alpha^2e_a^2}{Q^2y}
   [1+(1-y)^2]\notag\\
 &=\boxed{\frac{4\pi\alpha^2e_a^2}{Q^2y}(1-y)}.
\end{align}


\subsection{The factorized spin-dependent cross section}
\subsubsection{Collinear spin projections}

Use the same $\Phi\otimes\mathcal H\otimes\Delta$ structure as in Section 2.
Select different leading spin components:
\begin{equation}
 \begin{array}{c|c|c|c}
  \text{term}&\Phi\text{ projection}&\Delta\text{ projection}&chiral\\\hline
  UU&f_1&D_1&even\\
  UT&h_1&H_1^{\angle}&odd
 \end{array}
\end{equation}
\begin{itemize}
 \item $a$: flavor 
 \item Unpolarized parton distribution $f_1^a(x_B)$:  spin-averaged parton density.
 \item Transversity distribution $h_1^a(x_B)$: For a unit transversely polarized
 proton, $h_1^a(x_B)$= the difference of quark densities with transverse spins
 parallel and antiparallel to the proton spin.
 \item Unpolarized pair fragmentation function $D_1^a(z,c,M_h^2)$: describe $q \to h_1+h_2+X$.
 \item Interference fragmentation function $H_1^{\angle a}(z,c,M_h^2)$: 
 It correlates the outgoing quark spin with the pair orientation.
 \item Use the same scale $\mu\sim Q$ for all functions; suppress it below.
\end{itemize}
\begin{equation}
 \underbrace{h_1}_{\text{chiral odd}}
 \ \times\ \underbrace{H_1^{\angle}}_{\text{chiral odd}}.
\end{equation}
Initial spin $\to$ Quark spin $\to$ Hadron-pair angular distribution  


\subsubsection{Fraction matching}
At tree level, take the incoming quark momentum to be
$k=\xi\widetilde P$, where
$\widetilde P=P^+\bar n$ and $\widetilde P^2=0$.
The virtual photon strikes this quark:
\begin{equation}
 \gamma^*(q)+a(\xi\widetilde P)\longrightarrow a(k'),
 \qquad k'=\xi\widetilde P+q.
\end{equation}
\textbf{Put the outgoing quark on shell:} $k'^2=0$. Expand:
\begin{align}
 0=(\xi\widetilde P+q)^2
 &=\xi^2\widetilde P^2+2\xi\widetilde P\cdot q+q^2\notag\\
 &=2\xi\widetilde P\cdot q-Q^2\notag\\
 &\simeq 2\xi P\cdot q-Q^2\notag\\
 &=Q^2\left(\frac{\xi}{x_B}-1\right).
\end{align}
Here $\widetilde P^2=0$, $q^2=-Q^2$, and
$2P\cdot q=Q^2/x_B$. The replacement
$\widetilde P\cdot q\simeq P\cdot q$ keeps leading power.
Therefore $\boxed{\xi=x_B}$ at tree level.
For the pair use $P_h\simeq\zeta k'$ in the hard kinematics:
\begin{equation}
 z=\frac{P\cdot P_h}{P\cdot q}\simeq\zeta
 \quad\Rightarrow\quad\zeta=z.
\end{equation}


\subsubsection{The two structure functions}

At tree level the fraction convolutions collapse, giving
\begin{align}
 F_{UU}&=x_B\sum_{a=q,\bar q}e_a^2 f_1^a(x_B)D_1^a(z,c,M_h^2),\\
 F_{UT}^{\sin\Phi}
 &=\boxed{\frac{|\bm R_T|}{M_h}\,x_B
 \sum_{a=q,\bar q}e_a^2 h_1^a(x_B)H_1^{\angle a}(z,c,M_h^2)}.
 \label{eq:dihadron-structure}
\end{align}
\begin{itemize}
 \item $UU$: unpolarized electron and proton.
 \item $UT$: unpolarized electron and transversely polarized proton.
 \item Keep the \textbf{pair angle} even after integrating the total pair
 transverse momentum. This is why collinear transversity remains visible.
\end{itemize}

\subsubsection{The differential cross section}
Define the differential measure:
\begin{equation}
 d\Gamma=dx_B\,dy\,d\psi\,dz\,d\phi_R\,dM_h^2\,dc.
\end{equation}
Thus
\begin{equation}
 \frac{d\sigma}{d\Gamma}
 \equiv
 \frac{d\sigma}
 {dx_B\,dy\,d\psi\,dz\,d\phi_R\,dM_h^2\,dc}.
\end{equation}
\textbf{Find the hard weights.} From the unpolarized hard part in
Section 2 and the spin-transfer factor above,
\begin{align}
 Y_+&=1+(1-y)^2=2-2y+y^2,\\
 A(y)&=\frac{Y_+}{2}=1-y+\frac{y^2}{2},\\
 B(y)&=A(y)D_{NN}(y)
 =\frac{1+(1-y)^2}{2}
   \frac{2(1-y)}{1+(1-y)^2}
 =1-y.
\end{align}
\begin{equation}
 \boxed{\frac{d\sigma}{d\Gamma}
 =\frac{\alpha^2}{x_ByQ^2}
 \left[A(y)F_{UU}
 +|\bm S_T|B(y)\sin\Phi\,F_{UT}^{\sin\Phi}\right].}
 \label{eq:dihadron-cross-section}
\end{equation}
\begin{itemize}
 \item $A(y)F_{UU}$ survives when $|\bm S_T|=0$.
 \item $|\bm S_T|B(y)\sin\Phi F_{UT}^{\sin\Phi}$ changes sign when
 the target spin is reversed. Here $B(y)$ is the hard spin-transfer
 weight; $\sin\Phi$ measures the pair orientation relative to the spin.
\end{itemize}
\textbf{Insert the factorized structure functions}
from Eq.~\eqref{eq:dihadron-structure}. Suppress their arguments:
\begin{equation}
 \begin{aligned}
 \frac{d\sigma}{d\Gamma}
 =\frac{\alpha^2}{x_ByQ^2}\Bigg[
 &A(y)x_B\sum_a e_a^2f_1^aD_1^a\\
 &+|\bm S_T|B(y)\sin\Phi\,
   x_B\frac{|\bm R_T|}{M_h}
   \sum_a e_a^2h_1^aH_1^{\angle a}
 \Bigg].
 \end{aligned}
\end{equation}
Cancel the common $x_B$:
\begin{equation}
 \boxed{\begin{aligned}
 \frac{d\sigma}{d\Gamma}=\frac{\alpha^2}{Q^2y}
 \Bigg[&A(y)\sum_a e_a^2 f_1^aD_1^a\\
 &+|\bm S_T|B(y)\sin(\phi_R+\phi_S)
 \frac{|\bm R_T|}{M_h}\sum_a e_a^2h_1^aH_1^{\angle a}\Bigg].
 \end{aligned}}
\end{equation}
The products $f_1D_1$ and $h_1H_1^{\angle}$ show the
tree-level collinear factorization.
\textbf{Keep this angular measure and DiFF normalization together.}

\end{document}
